% GENERATED FILE. Edit canonical lessons, then run npm run generate:books.
% canonical-source-hash: fnv1a64:538d492d3bd2ffdf
% canonical-lessons: TE-C357-vamsam, TE-C357-taram, TE-C357-purvikulu, TE-C357-anatha, TE-C357-batuku, TE-R357-review-350-357

\chapter{Generations}
\label{ch:te-generations}
\hlchaptermodality{357}

\begin{chapteropening}
\textbf{By the end of this chapter:} \emph{I can say a lineage, a generation, ancestors, an orphan and to survive.}

Complete the last lesson of chapter 357: I can say a lineage, a generation, ancestors, an orphan and to survive.
\end{chapteropening}

\section[\texorpdfstring{\te{వంశం} (vaṁśaṁ)}{వంశం (vaṁśaṁ)}]{\te{వంశం} (vaṁśaṁ) — a lineage; a dynasty}

\label{lesson:TE-C357-vamsam}



\begin{quote}
Before the new one: say the Telugu for a farewell, then the Telugu for adoption.
\end{quote}

\subsection*{You'll want to know: \te{వంశం}}
\textbf{\te{వంశం}} — \emph{vaṁśaṁ} — \textquotedblleft{}a lineage; a dynasty\textquotedblright{}.

\begin{quote}
\te{ఆ} \te{రాజు} \te{గొప్ప} \te{వంశంలో} \te{పుట్టాడు}.
\end{quote}

\begin{quote}
\emph{ā rāju goppa vaṁśaṁlō puṭṭāḍu.}
\end{quote}

\begin{quote}
That king was born into a great dynasty.
\end{quote}

Plural: \te{వంశాలు}, \emph{vaṁśālu}.

\begin{cousinweb}
\textbf{\te{వంశం}} (\emph{vaṁśaṁ}) is borrowed from Sanskrit \emph{vaṁśa}, \textquotedblleft{}a bamboo\textquotedblright{}, and also \textquotedblleft{}a family, a lineage\textquotedblright{} (CDIAL 11175).
\end{cousinweb}

\begin{grammarlens}[title={What you've built}]
One more word for telling what happened.
\end{grammarlens}

\subsection*{Your turn}
\begin{itemize}
\raggedright
  \item \emph{Say it:} \emph{vaṁśaṁ}
  \item \emph{Say it:} \emph{vaṁśaṁ}, once more
  \item \emph{Say it:} \emph{ā rāju goppa vaṁśaṁlō puṭṭāḍu.} — That king was born into a great dynasty.
  \item \emph{Recall:} say the Telugu for a farewell, then the Telugu for adoption, then say \emph{vaṁśaṁ} again
\end{itemize}

\begin{tcolorbox}[breakable,colback=teal!4,colframe=teal!35!black,title={Before you move on}]
What does \te{వంశం} mean? (\textquotedblleft{}A lineage; a dynasty\textquotedblright{}.) Say it once more.
\end{tcolorbox}
\section[\texorpdfstring{\te{తరం} (taraṁ)}{తరం (taraṁ)}]{\te{తరం} (taraṁ) — a generation}

\label{lesson:TE-C357-taram}



\begin{quote}
Before the new one: say the Telugu for adoption, then the Telugu for a lineage.
\end{quote}

\subsection*{You'll want to know: \te{తరం}}
\textbf{\te{తరం}} — \emph{taraṁ} — \textquotedblleft{}a generation\textquotedblright{}.

\begin{quote}
\te{మా} \te{తరం} \te{పిల్లలు} \te{టీవీ} \te{చూస్తూ} \te{పెరిగారు}.
\end{quote}

\begin{quote}
\emph{mā taraṁ pillalu ṭīvī cūstū perigāru.}
\end{quote}

\begin{quote}
The children of our generation grew up watching TV.
\end{quote}

Plural: \te{తరాలు}, \emph{tarālu}.

\te{తరతరాలుగా}, for generations.

\begin{cousinweb}
\textbf{\te{తరం}} (\emph{taraṁ}) is the Telugu word for a generation; its origin is not settled.
\end{cousinweb}

\begin{grammarlens}[title={What you've built}]
One more word for telling what happened.
\end{grammarlens}

\subsection*{Your turn}
\begin{itemize}
\raggedright
  \item \emph{Say it:} \emph{taraṁ}
  \item \emph{Say it:} \emph{taraṁ}, once more
  \item \emph{Say it:} \emph{mā taraṁ pillalu ṭīvī cūstū perigāru.} — The children of our generation grew up watching TV.
  \item \emph{Recall:} say the Telugu for adoption, then the Telugu for a lineage, then say \emph{taraṁ} again
\end{itemize}

\begin{tcolorbox}[breakable,colback=teal!4,colframe=teal!35!black,title={Before you move on}]
What does \te{తరం} mean? (\textquotedblleft{}A generation\textquotedblright{}.) Say it once more.
\end{tcolorbox}
\section[\texorpdfstring{\te{పూర్వీకులు} (pūrvīkulu)}{పూర్వీకులు (pūrvīkulu)}]{\te{పూర్వీకులు} (pūrvīkulu) — ancestors; forefathers}

\label{lesson:TE-C357-purvikulu}



\begin{quote}
Before the new one: say the Telugu for a lineage, then the Telugu for a generation.
\end{quote}

\subsection*{You'll want to know: \te{పూర్వీకులు}}
\textbf{\te{పూర్వీకులు}} — \emph{pūrvīkulu} — \textquotedblleft{}ancestors; forefathers\textquotedblright{}.

\begin{quote}
\te{మా} \te{పూర్వీకులు} \te{రైతులు}.
\end{quote}

\begin{quote}
\emph{mā pūrvīkulu raitulu.}
\end{quote}

\begin{quote}
Our ancestors were farmers.
\end{quote}

Built on \te{పూర్వం}, formerly: those who came before.

\begin{cousinweb}
\textbf{\te{పూర్వీకులు}} (\emph{pūrvīkulu}) is built on Sanskrit \emph{pūrvika}, \textquotedblleft{}former, ancient\textquotedblright{} (Monier-Williams), from \emph{pūrva}, \textquotedblleft{}former\textquotedblright{} (CDIAL 8343).
\end{cousinweb}

\begin{grammarlens}[title={What you've built}]
One more word for telling what happened.
\end{grammarlens}

\subsection*{Your turn}
\begin{itemize}
\raggedright
  \item \emph{Say it:} \emph{pūrvīkulu}
  \item \emph{Say it:} \emph{pūrvīkulu}, once more
  \item \emph{Say it:} \emph{mā pūrvīkulu raitulu.} — Our ancestors were farmers.
  \item \emph{Recall:} say the Telugu for a lineage, then the Telugu for a generation, then say \emph{pūrvīkulu} again
\end{itemize}

\begin{tcolorbox}[breakable,colback=teal!4,colframe=teal!35!black,title={Before you move on}]
What does \te{పూర్వీకులు} mean? (\textquotedblleft{}Ancestors; forefathers\textquotedblright{}.) Say it once more.
\end{tcolorbox}
\section[\texorpdfstring{\te{అనాథ} (anātha)}{అనాథ (anātha)}]{\te{అనాథ} (anātha) — an orphan}

\label{lesson:TE-C357-anatha}



\begin{quote}
Before the new one: say the Telugu for a generation, then the Telugu for ancestors.
\end{quote}

\subsection*{You'll want to know: \te{అనాథ}}
\textbf{\te{అనాథ}} — \emph{anātha} — \textquotedblleft{}an orphan\textquotedblright{}.

\begin{quote}
\te{ఆ} \te{అనాథ} \te{పిల్లలను} \te{ఒక} \te{ఆశ్రమం} \te{చూసుకుంటుంది}.
\end{quote}

\begin{quote}
\emph{ā anātha pillalanu oka āśramaṁ cūsukuṇṭundi.}
\end{quote}

\begin{quote}
A home looks after those orphan children.
\end{quote}

Plural: \te{అనాథలు}, \emph{anāthalu}.

\begin{cousinweb}
\textbf{\te{అనాథ}} (\emph{anātha}) is borrowed from Sanskrit \emph{anātha}, \textquotedblleft{}having no master or protector\textquotedblright{} (Monier-Williams), from \emph{a-}, \textquotedblleft{}without\textquotedblright{}, and \emph{nātha}, \textquotedblleft{}a protector\textquotedblright{}.
\end{cousinweb}

\begin{grammarlens}[title={What you've built}]
One more word for telling what happened.
\end{grammarlens}

\subsection*{Your turn}
\begin{itemize}
\raggedright
  \item \emph{Say it:} \emph{anātha}
  \item \emph{Say it:} \emph{anātha}, once more
  \item \emph{Say it:} \emph{ā anātha pillalanu oka āśramaṁ cūsukuṇṭundi.} — A home looks after those orphan children.
  \item \emph{Recall:} say the Telugu for a generation, then the Telugu for ancestors, then say \emph{anātha} again
\end{itemize}

\begin{tcolorbox}[breakable,colback=teal!4,colframe=teal!35!black,title={Before you move on}]
What does \te{అనాథ} mean? (\textquotedblleft{}An orphan\textquotedblright{}.) Say it once more.
\end{tcolorbox}
\section[\texorpdfstring{\te{బతుకు} (batuku)}{బతుకు (batuku)}]{\te{బతుకు} (batuku) — to survive; to live}

\label{lesson:TE-C357-batuku}



\begin{quote}
Before the new one: say the Telugu for ancestors, then the Telugu for an orphan.
\end{quote}

\subsection*{You'll want to know: \te{బతుకు}}
\textbf{\te{బతుకు}} — \emph{batuku} — \textquotedblleft{}to survive; to live\textquotedblright{}.

\begin{quote}
\te{ఆ} \te{ప్రమాదంలో} \te{అతను} \te{బతికాడు}.
\end{quote}

\begin{quote}
\emph{ā pramādaṁlō atanu batikāḍu.}
\end{quote}

\begin{quote}
He survived that accident.
\end{quote}

Past (I …): \te{బతికాను}, \emph{batikānu}.

Also a noun: \te{బతుకు}, a life, a living.

\begin{cousinweb}
\textbf{\te{బతుకు}} (\emph{batuku}) is native Dravidian: DEDR 5372 lists Telugu \emph{batuku}, \emph{braduku}, \textquotedblleft{}to live, subsist, survive\textquotedblright{}.
\end{cousinweb}

\begin{grammarlens}[title={What you've built}]
One more word for telling what happened.
\end{grammarlens}

\subsection*{Your turn}
\begin{itemize}
\raggedright
  \item \emph{Say it:} \emph{batuku}
  \item \emph{Say it:} \emph{batuku}, once more
  \item \emph{Say it:} \emph{ā pramādaṁlō atanu batikāḍu.} — He survived that accident.
  \item \emph{Recall:} say the Telugu for ancestors, then the Telugu for an orphan, then say \emph{batuku} again
\end{itemize}

\begin{tcolorbox}[breakable,colback=teal!4,colframe=teal!35!black,title={Before you move on}]
What does \te{బతుకు} mean? (\textquotedblleft{}To survive; to live\textquotedblright{}.) Say it once more.
\end{tcolorbox}
\section[(dialogue)]{Words from chapters 350–357 — a review}

\label{lesson:TE-R357-review-350-357}



\begin{quote}
Say the Telugu for a crime and a criminal.
\end{quote}

\subsection*{Your turn}
Say these aloud:

\begin{itemize}
\raggedright
  \item a crime, a criminal, a guilty person, an innocent person, to punish
  \item a captive, to detain, to capture, to release, to free from captivity
  \item to pursue, to hide oneself, to surrender, to turn in, to step aside
  \item a disappearance, to vanish, to fade from view, to appear suddenly, to detect
  \item a conspiracy, to mislead, to expose, to cover up, a spy
  \item an anniversary, a wedding anniversary, a sixtieth birthday, a birth anniversary, a felicitation
  \item a public celebration, to felicitate, a housewarming, a farewell, adoption
  \item a lineage, a generation, ancestors, an orphan, to survive
\end{itemize}

\begin{tcolorbox}[breakable,colback=teal!4,colframe=teal!35!black,title={Before you move on}]
A crime? (\textbf{\te{నేరం}}, \emph{nēraṁ}.) A criminal? (\textbf{\te{నేరస్థుడు}}, \emph{nērasthuḍu}.) A guilty person? (\textbf{\te{దోషి}}, \emph{dōṣi}.) An innocent person? (\textbf{\te{నిర్దోషి}}, \emph{nirdōṣi}.) To punish? (\textbf{\te{శిక్షించు}}, \emph{śikṣiñcu}.) A captive? (\textbf{\te{బందీ}}, \emph{bandī}.) To detain? (\textbf{\te{నిర్బంధించు}}, \emph{nirbandhiñcu}.) To capture? (\textbf{\te{బంధించు}}, \emph{bandhiñcu}.) To release? (\textbf{\te{విడుదల} \te{చేయు}}, \emph{viḍudala cēyu}.) To free from captivity? (\textbf{\te{విడిపించు}}, \emph{viḍipiñcu}.) To pursue? (\textbf{\te{వెంబడించు}}, \emph{vembaḍiñcu}.) To hide oneself? (\textbf{\te{దాక్కో}}, \emph{dākkō}.) To surrender? (\textbf{\te{లొంగిపోవు}}, \emph{loṅgipōvu}.) To turn in? (\textbf{\te{పట్టించు}}, \emph{paṭṭiñcu}.) To step aside? (\textbf{\te{తప్పుకో}}, \emph{tappukō}.) A disappearance? (\textbf{\te{అదృశ్యం}}, \emph{adṛśyaṁ}.) To vanish? (\textbf{\te{మాయమవు}}, \emph{māyamavu}.) To fade from view? (\textbf{\te{కనుమరుగవు}}, \emph{kanumarugavu}.) To appear suddenly? (\textbf{\te{ప్రత్యక్షమవు}}, \emph{pratyakṣamavu}.) To detect? (\textbf{\te{కనిపెట్టు}}, \emph{kanipeṭṭu}.) A conspiracy? (\textbf{\te{కుట్ర}}, \emph{kuṭra}.) To mislead? (\textbf{\te{మభ్యపెట్టు}}, \emph{mabhyapeṭṭu}.) To expose? (\textbf{\te{బయటపెట్టు}}, \emph{bayaṭapeṭṭu}.) To cover up? (\textbf{\te{కప్పిపుచ్చు}}, \emph{kappipuccu}.) A spy? (\textbf{\te{గూఢచారి}}, \emph{gūḍhacāri}.) An anniversary? (\textbf{\te{వార్షికోత్సవం}}, \emph{vārṣikōtsavaṁ}.) A wedding anniversary? (\textbf{\te{పెళ్ళిరోజు}}, \emph{peḷḷirōju}.) A sixtieth birthday? (\textbf{\te{షష్టిపూర్తి}}, \emph{ṣaṣṭipūrti}.) A birth anniversary? (\textbf{\te{జయంతి}}, \emph{jayanti}.) A felicitation? (\textbf{\te{సన్మానం}}, \emph{sanmānaṁ}.) A public celebration? (\textbf{\te{ఉత్సవం}}, \emph{utsavaṁ}.) To felicitate? (\textbf{\te{సన్మానించు}}, \emph{sanmāniñcu}.) A housewarming? (\textbf{\te{గృహప్రవేశం}}, \emph{gṛhapravēśaṁ}.) A farewell? (\textbf{\te{వీడ్కోలు}}, \emph{vīḍkōlu}.) Adoption? (\textbf{\te{దత్తత}}, \emph{dattata}.) A lineage? (\textbf{\te{వంశం}}, \emph{vaṁśaṁ}.) A generation? (\textbf{\te{తరం}}, \emph{taraṁ}.) Ancestors? (\textbf{\te{పూర్వీకులు}}, \emph{pūrvīkulu}.) An orphan? (\textbf{\te{అనాథ}}, \emph{anātha}.) To survive? (\textbf{\te{బతుకు}}, \emph{batuku}.)
\end{tcolorbox}
